% Capítulo 0: estrategias de aprendizaje
\section*{Objetivos del curso}
\begin{itemize}[leftmargin=*]
	\item Alcanzar un nivel A1 sólido en comprensión y producción oral y escrita.
	\item Construir hábitos diarios de estudio (15--30 minutos) con enfoque en input comprensible y producción guiada.
	\item Preparar un cuaderno de trabajo y un espacio digital (Anki/flashcards) para repaso espaciado.
\end{itemize}

\section*{Rutina diaria sugerida}
\begin{enumerate}[leftmargin=*]
	\item 5--10 minutos: repaso con repetición espaciada (palabras y frases de días previos).
	\item 10--15 minutos: lectura o escucha de material corto con apoyo bilingüe.
	\item 5--10 minutos: práctica activa (dictados, traducción inversa, shadowing, grabarte hablando).
\end{enumerate}

\section*{Técnicas clave}
\begin{itemize}[leftmargin=*]
	\item \textbf{Input comprensible}: escoger textos y audios al 90\% de comprensión, subrayar estructuras clave.
	\item \textbf{Shadowing}: repetir en voz alta inmediatamente después de oír, priorizando ritmo y entonación.
	\item \textbf{Bloques léxicos}: aprender frases completas en vez de palabras aisladas.
	\item \textbf{Espaciado}: programar repasos 1-3-7-14 días después de aprender algo nuevo.
\end{itemize}

\section*{Recursos recomendados}
\begin{itemize}[leftmargin=*]
	\item Diccionarios: Van Dale en línea, Glosbe (ejemplos), Wiktionary (pronunciación IPA).
	\item Escucha: podcasts A1, diálogos cortos en YouTube, TTS neerlandés para practicar ritmo.
	\item Escritura: diarios breves; pedir correcciones con lenguaje sencillo.
\end{itemize}

\section*{Seguimiento y evaluación}
\begin{itemize}[leftmargin=*]
	\item Establece micro-metas semanales (e.g., ``presentarme en 6 frases claras'').
	\item Usa grabaciones mensuales para comparar progreso en pronunciación.
	\item Autoevaluación con checklists A1 del MCER al finalizar cada bloque de 5 capítulos.
\end{itemize}
